% Bewertungspaket zum Gesamttest «Bruchrechnen» — Quelle des PDFs (python3 scripts/build-lp-pdf.py bruchrechnen).
% Ganze Punkte je Teilschritt; Werte und Fehlerwerte mit python3 nachgerechnet (scripts/lp/bruchrechnen/zahlen.py, 09.10.2026).
% Nach der Prüfung vom 09.10.2026 überarbeitet: neue Aufgaben G1(b), G3, G4, G5(Skizze), G6, G7; einheitliche Restpunkte;
% Dezimal- und gemischte Zahlen unter «Allgemein»; Testblatt mit hochladen.
\documentclass[11pt]{article}
\usepackage{../lp-druck}
\renewcommand{\lpname}{Bruchrechnen}
\renewcommand{\lpteilgebiet}{Arithmetik/Algebra ·}
\renewcommand{\lpfuss}{Bewertungspaket}
\newcolumntype{P}{>{\centering\arraybackslash\rule[-2.6ex]{0pt}{6.2ex}}p{10mm}}  % Stütze: Platz für \dfrac oben und unten
% Hohe Brüche (\dfrac) brauchen Luft zwischen den Zeilen, sonst berühren sie die Nachbarzeile (Sichtprüfung 09.10.2026).
\newenvironment{raster}{\par\smallskip\noindent\renewcommand{\arraystretch}{1.15}\linespread{1.35}\selectfont\tabularx{\linewidth}{@{}>{\raggedright\arraybackslash}p{38mm}P X@{}}\toprule
  \small\textsc{Teilschritt} & \small\textsc{P} & \small\textsc{Musterlösung}\\\midrule}{\bottomrule\endtabularx\par}
\newcommand{\fehler}[1]{\par\smallskip{\small\textcolor{lprot}{\bfseries Typische Fehler:} #1}}
\newcommand{\E}{\,{\footnotesize\bfseries\color{lpgruen}(E)}}
\begin{document}
\lpkopf{Bewertungspaket zum Gesamttest}{}

\begin{lpachtung}{Erst lösen, dann öffnen}
Dieses Paket enthält die Musterlösung. Wer es vor dem Lösen liest, misst am Ende nur, wie gut das Abschreiben gelingt.
\end{lpachtung}

\section*{1 · So gehst du vor}
\begin{enumerate}[leftmargin=*,itemsep=2pt]
\item \textbf{Gesamttest lösen} — vollständig, mit Rechenweg, ohne dieses Paket und ohne Taschenrechner.
\item \textbf{Selbst bewerten} mit Abschnitt~3: Zeile für Zeile vergleichen und die Punkte eintragen. Das geht ohne Hilfsmittel
  und ist der Normalfall.
\item \textbf{Oder bewerten lassen.} Ob du dafür eine KI nutzt und welche, entscheidest du selbst und in eigener Verantwortung —
  je nachdem, was dir aufgrund deines Alters und deiner persönlichen Situation erlaubt ist (Altersgrenzen und
  Nutzungsbedingungen des Anbieters, allenfalls Einverständnis deiner Eltern). Fotografiere oder scanne deine Lösung gut lesbar,
  \textbf{ohne Namen}. Handyfotos tragen oft unsichtbar Standort, Zeit und Gerät mit (Metadaten) — lade lieber einen Scan oder ein
  Bildschirmfoto des Fotos hoch. Lade nur deine Lösung, das Testblatt und dieses PDF hoch, keine persönlichen Daten, und füge den Auftrag aus
  Abschnitt~2 ein.
\item \textbf{Rückmeldung prüfen.} Stimmt, was die KI aus deiner Handschrift gelesen hat? Eine KI kann sich irren —
  im Zweifel gilt das Raster in Abschnitt~3.
\item \textbf{Gezielt wiederholen} nach Abschnitt~4.
\end{enumerate}

\section*{2 · Auftrag an die KI}
\begin{tcolorbox}[colback=lplinie!25,colframe=lplinie,boxrule=0.5pt,arc=1mm,fontupper=\small]
Du bist Korrektorin oder Korrektor für Mathematik auf der Sekundarstufe I in der Schweiz. Im Anhang findest du (1)~das Testblatt mit den Aufgaben, (2)~das Bewertungspaket zum Gesamttest «Bruchrechnen» mit Musterlösung und Punkteraster und (3)~Fotos meiner handschriftlichen Lösung.\par\medskip
Geh so vor:\par
1. Schreib zu jeder Aufgabe G1 bis G9 zuerst ab, was du in meiner Lösung liest — Rechenweg und Ergebnis. Was du nicht sicher lesen kannst, markierst du mit [unsicher]. Nicht raten. Achte besonders auf Bruchstriche und Minuszeichen.\par
2. Vergib die Punkte genau nach dem Raster im Bewertungspaket (Abschnitt 3), ganze Punkte je Zeile. Ziehe einen Fehler nur einmal ab: Rechne ich mit einem falschen Zwischenergebnis richtig weiter, gibt es die Folgepunkte. Ausnahme: Zeilen mit \textbf{(E)} gibt es nur für das richtige Ergebnis, nie als Folgepunkt. Die Zeilen «Typische Fehler» sagen, welche Rasterzeile bei diesem Fehler 0 Punkte gibt — sie sind kein zusätzlicher Abzug.\par
3. Andere richtige Lösungswege zählen voll, zum Beispiel ein anderer gemeinsamer Nenner als der Hauptnenner oder Kürzen erst am Schluss. Ein Ergebnis zählt nur \textbf{vollständig gekürzt}: Ein richtiger, aber ungekürzter Bruch gibt in der (E)-Zeile 0 Punkte. Gleichwertig: das Minus vor dem Bruch, im Zähler oder im Nenner; eine ganze Zahl als $1$ oder $\tfrac{1}{1}$. Für Dezimalzahlen und gemischte Zahlen gilt die Regel «Schreibweise des Ergebnisses» in Abschnitt~3.\par
4. Begründe jeden Abzug in einem Satz und nenne die Fehlerart (zum Beispiel «Nenner addiert», «Kehrwert vergessen», «Punkt vor Strich»). Schreib die Musterlösung nicht ab — zeig mir, wo mein Fehler liegt.\par
5. Gib zum Schluss eine Tabelle «Aufgabe | Punkte | Begründung», die Gesamtpunktzahl (von 25) und — nach der Selbsteinschätzung in Abschnitt 4 — welches Kapitel des Leitprogramms ich wiederholen soll. Keine Note.\par
6. Fehlt ein Foto oder ist eine Aufgabe unlesbar, sag es, statt sie zu bewerten.
\end{tcolorbox}

\needspace{16\baselineskip}
\section*{3 · Musterlösung und Punkteraster}
\textbf{Allgemein:} Ganze Punkte je Zeile. Ein Fehler wird nur einmal abgezogen (Folgefehler). Gleichwertige Wege zählen voll.
In der Spalte P: \textbf{(E)}~= Ergebnispunkt, Ergebnis richtig, vollständig gekürzt \emph{und} Weg erkennbar; ohne Zeichen = die
Zeile bewertet den Weg oder die Begründung. Der ganze Test ist ohne Taschenrechner gestellt.
«Typische Fehler» nennt, welche Zeile 0 Punkte gibt, und die Punkte, die bleiben.
\textbf{Teilrichtige Zeilen:} Verlangt eine Zeile mehrere Angaben und stimmt nur ein Teil, gibt die Zeile 0 Punkte —
ausser die Zeile nennt ausdrücklich eine Aufteilung.
\textbf{Schreibweise des Ergebnisses:} Verlangt ist ein vollständig gekürzter Bruch oder eine ganze Zahl. Eine \emph{exakte}
Dezimalzahl ($-0.5$ statt $-\tfrac{1}{2}$) oder eine gemischte Zahl ($2\tfrac{1}{3}$ statt $\tfrac{7}{3}$) gibt den
Ergebnispunkt trotzdem, mit dem Hinweis, dass ein Bruch verlangt war. Eine gerundete oder abgebrochene Dezimalzahl
($-2.33$ für $-\tfrac{7}{3}$) gibt 0 — sie ist nicht gleich dem Ergebnis. Ein ungekürzter Bruch gibt im Ergebnispunkt 0.

\begin{lpaufgabe}{G1}{Kürzen}{2 P}
\begin{raster}
(a) & 1\E & $\text{ggT}(84;\ 126) = 42$: $\tfrac{84}{126} = \tfrac{2}{3}$. Ein Kürzschritt oder der ggT muss dastehen; in mehreren Schritten gekürzt ist gleichwertig, solange am Ende $\tfrac{2}{3}$ steht.\\
(b) & 1 & Ja: Beide lassen sich auf $\tfrac{3}{4}$ kürzen ($21 : 7$, $28 : 7$ und $15 : 5$, $20 : 5$). Gleichwertig: Erweitern auf denselben Nenner oder $21 \cdot 20 = 420 = 28 \cdot 15$. Nur «ja» ohne Begründung: 0.\\
\end{raster}
\fehler{(a) Nur teilweise gekürzt ($\tfrac{42}{63}$, $\tfrac{14}{21}$, $\tfrac{12}{18}$): Zeile (a) 0 $\to$ höchstens 1 von 2\,P.}
\end{lpaufgabe}

\begin{lpaufgabe}{G2}{Drei Brüche}{3 P}
\begin{raster}
Hauptnenner & 1 & $\text{kgV}(15;\ 6;\ 2) = 30$. Jeder gemeinsame Nenner zählt (z.\,B. 180), wenn er richtig verwendet wird.\\
Erweitern & 1 & $\dfrac{14}{30} + \dfrac{25}{30} - \dfrac{15}{30}$\\
Ergebnis & 1\E & $\dfrac{24}{30} = \dfrac{4}{5}$\\
\end{raster}
\fehler{Zähler und Nenner je verrechnet ($\tfrac{7 + 5 - 1}{15 + 6 - 2} = \tfrac{11}{19}$): alle Zeilen 0.
Richtig bis $\tfrac{24}{30}$, nicht gekürzt: Ergebniszeile 0 $\to$ 2 von 3\,P.
Nur den Nenner erweitert, Zähler gelassen ($\tfrac{7}{30} + \tfrac{5}{30} - \tfrac{1}{30} = \tfrac{11}{30}$): Erweitern- und Ergebniszeile 0 $\to$ 1 von 3\,P.}
\end{lpaufgabe}

\begin{lpaufgabe}{G3}{Vergleichen und Unterschied}{2 P}
\begin{raster}
Gleichnamig & 1 & Hauptnenner 36: $\tfrac{3}{4} = \tfrac{27}{36}$, $\tfrac{7}{9} = \tfrac{28}{36}$. Jeder andere gemeinsame Nenner zählt ebenso.\\
Ergebnis & 1\E & $\tfrac{7}{9}$ ist grösser, um $\tfrac{28}{36} - \tfrac{27}{36} = \tfrac{1}{36}$. Beide Angaben nötig.\\
\end{raster}
\fehler{«$\tfrac{3}{4}$ ist grösser» (z.\,B. weil der Nenner kleiner ist): Ergebniszeile 0; ist richtig gleichnamig gemacht $\to$ 1 von 2\,P.
Unterschied als $\tfrac{7 - 3}{9 - 4} = \tfrac{4}{5}$: beide Zeilen 0.
Richtig verglichen, Unterschied fehlt oder falsch: Ergebniszeile 0 $\to$ 1 von 2\,P.}
\end{lpaufgabe}

\begin{lpaufgabe}{G4}{Mal und durch}{3 P}
\begin{raster}
Multiplizieren & 1 & $\dfrac{9}{14} \cdot \dfrac{7}{12} = \dfrac{63}{168} = \dfrac{3}{8}$ (oder vorher kürzen: $7$ mit $14$, $9$ mit $12$)\\
Kehrwert & 1 & $\dfrac{3}{8} : \dfrac{3}{4} = \dfrac{3}{8} \cdot \dfrac{4}{3}$ — Kehrwert des zweiten Bruchs\\
Ergebnis & 1\E & $\dfrac{12}{24} = \dfrac{1}{2}$\\
\end{raster}
\fehler{Kehrwert vergessen ($\tfrac{3}{8} \cdot \tfrac{3}{4} = \tfrac{9}{32}$): Kehrwert- und Ergebniszeile 0 $\to$ 1 von 3\,P.
Kehrwert vom falschen Bruch ($\tfrac{8}{3} \cdot \tfrac{3}{4} = 2$): Kehrwert- und Ergebniszeile 0 $\to$ 1 von 3\,P.
Zuerst $\tfrac{7}{12} : \tfrac{3}{4} = \tfrac{7}{9}$ gerechnet, dann $\tfrac{9}{14} \cdot \tfrac{7}{9} = \tfrac{1}{2}$: hier gleichwertig, weil
$a \cdot b : c = a \cdot (b : c)$ gilt — volle Punkte. (Umgekehrt, bei $a : b \cdot c$, gilt das nicht: dort von links nach rechts.)}
\end{lpaufgabe}

\begin{lpaufgabe}{G5}{Flasche und Gläser}{3 P}
\begin{raster}
(a) Rechnung & 1\E & $\tfrac{3}{2} : \tfrac{1}{4} = \tfrac{3}{2} \cdot 4 = 6$ Gläser. Nur «6» ohne Rechnung: 0.\\
(a) Skizze & 1 & Ein Streifen für $\tfrac{3}{2}$ Liter (anderthalb Ganze), eingeteilt in Viertel: sechs Stücke zu $\tfrac{1}{4}$ sichtbar. Auch ein Streifen von 1 Liter in Vierteln plus ein halber Streifen zählt. Ungenau gezeichnet ist erlaubt, die Einteilung muss stimmen.\\
(b) & 1\E & $\tfrac{2}{3} \cdot \tfrac{3}{2} = 1$ (Liter). Die Einheit ist nicht verlangt.\\
\end{raster}
\fehler{(a) $\tfrac{3}{2} \cdot \tfrac{1}{4} = \tfrac{3}{8}$ (mal statt durch): Rechnungszeile 0; eine richtige Skizze gibt trotzdem ihren Punkt.
(b) $\tfrac{3}{2} : \tfrac{2}{3} = \tfrac{9}{4}$ («von» als durch gelesen): Zeile (b) 0.}
\end{lpaufgabe}

\begin{lpaufgabe}{G6}{Klammer und Punkt vor Strich}{3 P}
\begin{raster}
Klammer & 1 & $\dfrac{5}{6} - \dfrac{1}{3} = \dfrac{5}{6} - \dfrac{2}{6} = \dfrac{3}{6} = \dfrac{1}{2}$\\
Punkt vor Strich & 1 & $\dfrac{1}{2} \cdot \dfrac{3}{4} = \dfrac{3}{8}$ — vor der Addition\\
Ergebnis & 1\E & $\dfrac{3}{8} + \dfrac{1}{8} = \dfrac{4}{8} = \dfrac{1}{2}$\\
\end{raster}
\fehler{Zuerst addiert, $\tfrac{1}{2} \cdot \left(\tfrac{3}{4} + \tfrac{1}{8}\right) = \tfrac{7}{16}$: Punkt-vor-Strich- und Ergebniszeile 0 $\to$ 1 von 3\,P.
Klammer weggelassen, $\tfrac{5}{6} - \tfrac{1}{3} \cdot \tfrac{3}{4} + \tfrac{1}{8} = \tfrac{17}{24}$: Klammer- und Ergebniszeile 0, Punkt vor Strich richtig angewandt (zuerst $\tfrac{1}{3} \cdot \tfrac{3}{4}$) gibt die mittlere Zeile $\to$ 1 von 3\,P.
Ergebnis $\tfrac{4}{8}$ nicht gekürzt: Ergebniszeile 0 $\to$ 2 von 3\,P.}
\end{lpaufgabe}

\begin{lpaufgabe}{G7}{Potenz}{3 P}
\begin{raster}
Potenz & 1 & $\left(-\dfrac{3}{4}\right)^2 = +\dfrac{9}{16}$\\
Division & 1 & $\dfrac{9}{16} : \dfrac{9}{8} = \dfrac{9}{16} \cdot \dfrac{8}{9} = \dfrac{1}{2}$ — vor der Subtraktion\\
Ergebnis & 1\E & $\dfrac{1}{2} - 1 = -\dfrac{1}{2}$\\
\end{raster}
\fehler{Quadrat negativ ($-\tfrac{9}{16}$, dann $-\tfrac{1}{2} - 1 = -\tfrac{3}{2}$): Potenzzeile 0, Division als Folgefehler 1, Ergebniszeile 0 $\to$ 1 von 3\,P.
Nur den Zähler quadriert ($\tfrac{9}{4}$, dann $\tfrac{9}{4} : \tfrac{9}{8} = 2$, Ergebnis $1$): Potenzzeile 0, Division als Folgefehler 1, Ergebniszeile 0 $\to$ 1 von 3\,P.
Kehrwert vom falschen Bruch ($\tfrac{16}{9} \cdot \tfrac{9}{8} = 2$, Ergebnis $1$): Divisions- und Ergebniszeile 0 $\to$ 1 von 3\,P.
Zuerst $\tfrac{9}{8} - 1$ gerechnet ($\tfrac{9}{16} : \tfrac{1}{8} = \tfrac{9}{2}$): Divisions- und Ergebniszeile 0 $\to$ 1 von 3\,P.}
\end{lpaufgabe}

\begin{lpaufgabe}{G8}{Negative Brüche}{3 P}
\begin{raster}
Kehrwert & 1 & $\dfrac{2}{3} : \left(-\dfrac{4}{9}\right) = \dfrac{2}{3} \cdot \left(-\dfrac{9}{4}\right)$ — Kehrwert des zweiten Bruchs\\
Vorzeichen und Wert & 1 & $= -\dfrac{18}{12} = -\dfrac{3}{2}$ (Plus mal Minus gibt Minus)\\
Ergebnis & 1\E & $-\dfrac{5}{6} - \dfrac{3}{2} = -\dfrac{5}{6} - \dfrac{9}{6} = -\dfrac{14}{6} = -\dfrac{7}{3}$\\
\end{raster}
\fehler{Vorzeichen der Division falsch ($+\tfrac{3}{2}$, dann $-\tfrac{5}{6} + \tfrac{3}{2} = \tfrac{2}{3}$): Vorzeichenzeile 0, Ergebniszeile 0 $\to$ 1 von 3\,P.
Von links nach rechts ($\left(-\tfrac{5}{6} + \tfrac{2}{3}\right) : \left(-\tfrac{4}{9}\right) = \tfrac{3}{8}$): Kehrwert richtig gebildet gibt die erste Zeile; Vorzeichen- und Ergebniszeile 0 $\to$ 1 von 3\,P (wie G6 und G7).
Kehrwert vom falschen Bruch ($\tfrac{3}{2} \cdot \left(-\tfrac{4}{9}\right) = -\tfrac{2}{3}$, dann $-\tfrac{3}{2}$): Kehrwertzeile 0, Vorzeichen als Folgefehler 1, Ergebniszeile 0 $\to$ 1 von 3\,P.
Ergebnis $-\tfrac{14}{6}$ nicht gekürzt: Ergebniszeile 0 $\to$ 2 von 3\,P.}
\end{lpaufgabe}

\begin{lpaufgabe}{G9}{Fehler finden}{3 P}
\begin{raster}
Regel & 1 & \emph{Punkt vor Strich} verletzt: $\tfrac{1}{6} \cdot 4$ muss vor der Addition gerechnet werden; die Lösung hat zuerst addiert. Nur «falsch gerechnet» ohne Regel: 0.\\
Produkt & 1 & $\dfrac{1}{6} \cdot 4 = \dfrac{4}{6} = \dfrac{2}{3}$\\
Ergebnis & 1\E & $\dfrac{1}{3} + \dfrac{2}{3} = 1$\\
\end{raster}
\fehler{Regel richtig genannt, aber $\tfrac{1}{6} \cdot 4 = \tfrac{4}{24}$ (auch den Nenner multipliziert): Produktzeile 0; $\tfrac{1}{3} + \tfrac{1}{6} = \tfrac{1}{2}$ als Folgefehler ist kein Ergebnispunkt $\to$ 1 von 3\,P.
Fehler bei $\tfrac{1}{3} + \tfrac{1}{6} = \tfrac{3}{6}$ gesucht (dieser Schritt ist richtig): Regelzeile 0; Produkt und Ergebnis richtig neu gerechnet $\to$ 2 von 3\,P.}
\end{lpaufgabe}

\needspace{12\baselineskip}
\section*{4 · Selbsteinschätzung}
\begin{tabularx}{\linewidth}{@{}l X@{}}
\toprule
\textbf{22 – 25 P} & Die geprüften Teile sitzen. Wo du Punkte verloren hast: das Kapitel dieser Aufgabe nochmals.\\
\textbf{17 – 21 P} & Den schwächsten Teil nochmals — gemessen am Anteil (Teil A 7, B 6, C 12 P): Tüfteln und Übungen der Kapitel dieses Teils.\\
\textbf{11 – 16 P} & Zurück zu den Kapiteln aller Aufgaben, in denen du Punkte verloren hast.\\
\textbf{0 – 10 P} & Zurück zu den Kapiteln jedes Teils, in dem du weniger als die Hälfte erreicht hast; sind es alle drei, zum Vortest und zu Kapitel 1.\\
\bottomrule
\end{tabularx}

\smallskip
\textbf{Aufgabe $\to$ Kapitel:} G1 $\to$ 1; G2 $\to$ 2; G3 $\to$ 1 und 2; G4, G5 $\to$ 3; G6 bis G9 $\to$ 4.
\end{document}
